% domin2.tex（与 domin1 行列参数对齐）
\begin{tikzpicture}[
    scale=0.42,
    every node/.style={font=\small},
    gridline/.style={gray!60, very thin},
    proclabel/.style={font=\normalsize\bfseries, text=brown!70!black},
    arrowstyle/.style={-{Stealth[length=2.5mm]}, thick, purple!70!black}
]

\def\r{3}        % 辅助层列数
\def\NxL{4}      % 内部单元列数（与 domin1 单进程宽度一致）
\def\NyL{9}      % 行数（与 domin1 完全一致）

% 辅助网格层（灰色）
\fill[gray!40, opacity=0.85] (0,0)        rectangle (\r,       \NyL);
\fill[gray!40, opacity=0.85] (\r+\NxL,0)  rectangle (2*\r+\NxL,\NyL);

% 内部单元（蓝色，与 domin1 进程2 完全一致）
\fill[blue!60, opacity=0.85] (\r,0) rectangle (\r+\NxL, \NyL);

% 网格线
\foreach \x in {0,...,10} {
    \draw[gridline] (\x,0) -- (\x,\NyL);
}
\foreach \y in {0,...,\NyL} {
    \draw[gridline] (0,\y) -- (2*\r+\NxL,\y);
}

% 分隔粗线
\draw[very thick, black] (\r,0)      -- (\r,\NyL);
\draw[very thick, black] (\r+\NxL,0) -- (\r+\NxL,\NyL);
% 外框
\draw[very thick, black] (0,0) rectangle (2*\r+\NxL, \NyL);

% 标注
\node[proclabel] at (\r+0.5*\NxL, 0.5*\NyL) {\small 本进程};

% 通信箭头
\draw[arrowstyle] (-1.8, \NyL/2+0.8) -- (0.4, \NyL/2+0.8)
    node[midway, above, font=\footnotesize, text=black] {接收};
\draw[arrowstyle] (\r+0.4, \NyL/2-0.8) -- (-1.8, \NyL/2-0.8)
    node[midway, above, font=\footnotesize, text=black] {发送};
\draw[arrowstyle] (\r+\NxL-0.4, \NyL/2+0.8) -- (2*\r+\NxL+1.8, \NyL/2+0.8)
    node[midway, above, font=\footnotesize, text=black] {发送};
\draw[arrowstyle] (2*\r+\NxL+1.8, \NyL/2-0.8) -- (\r+\NxL+0.4, \NyL/2-0.8)
    node[midway, above, font=\footnotesize, text=black] {接收};


\end{tikzpicture}